\begin{lemma}[Peripheral isolation of pulled full continua]
\label{lem:peripheral-isolation}
Let a chain of \(p\) quadratic-like pullbacks be given (a composition
of proper degree-two holomorphic maps between plane topological
disks), and let \(\mathcal B\) be a finite pairwise-disjoint family of
full continua compactly contained in the terminal disk.
\begin{enumerate}
\item \emph{(Finite full pullback.)}  The inverse image of one full
      continuum under the chain is a union of at most \(2^p\) pairwise
      disjoint full continua; every resulting component maps onto its
      parent continuum with positive integer degree.  The inverse image
      of \(\mathcal B\) is again a finite pairwise-disjoint family of
      full continua.
\item \emph{(Isolating Jordan disks.)}  If \(C_1,\dots,C_N\Subset U\)
      are pairwise disjoint full continua in a plane topological disk
      \(U\), then for each \(i\) there is a Jordan domain
      \(D_i\Subset U\) with
      \[
       C_i\subset D_i,\qquad
       \overline D_i\cap C_j=\varnothing\quad(j\ne i),
      \]
      and the \(D_i\) may be chosen pairwise disjoint.
\item \emph{(Peripheral connectivity.)}  With \(C_1,\dots,C_N\) as in
      (2), the complement \(U\setminus\bigcup_{i=1}^N C_i\) is path
      connected.
\end{enumerate}
\end{lemma}

\begin{proof}
(1)  Apply Lemma~\ref{lem:quadratic-pullback-dichotomy} successively
along the chain.  At each step, each parent continuum has either one
degree-two child or two degree-one children, and every child is mapped
onto its parent; children of different parents are disjoint because
their images lie in disjoint parent continua.  Induction over the
\(p\) steps gives the bound \(2^p\), the stated surjectivity, and a
positive product degree for every terminal component.  Applying this
separately to the members of the finite family \(\mathcal B\) proves
the final statement of (1).

(2)  The finite family \(C_1,\dots,C_N\) has positive minimum pairwise
distance and positive distance from \(\partial U\).  Apply
Lemma~\ref{lem:jordan-basis} inside pairwise-disjoint metric
neighborhoods of the \(C_i\).

(3)  Applying Lemma~\ref{lem:jordan-basis} twice inside the disjoint
neighborhoods used in (2), choose pairwise-disjoint Jordan domains
\[
 C_i\subset D_i^-,\qquad
 \overline{D_i^-}\subset D_i^+,\qquad
 \overline{D_i^+}\subset U .
\]
The finitely connected plane domain
\[
 U_0=U\setminus\bigcup_{i=1}^N\overline{D_i^-}
\]
is connected: after a homeomorphism \(U\cong\mathbb D\), removing
finitely many pairwise-disjoint closed Jordan disks compactly
contained in \(\mathbb D\) leaves a disk with finitely many holes.  As
an open connected subset of the plane, \(U_0\) is path connected.

For each \(i\), the open set \(D_i^+\setminus C_i\) is connected.
Indeed, \(\mathbb C\setminus C_i\) is connected by fullness, and its
intersection with the Jordan disk \(D_i^+\) cannot have a component
separated from the outer collar of \(\partial D_i^+\): such a
component would be a bounded complementary component of \(C_i\),
contradicting fullness.  Hence \(D_i^+\setminus C_i\) is path
connected.  It contains the nonempty annulus
\(D_i^+\setminus\overline{D_i^-}\), which also lies in \(U_0\); thus
every \(D_i^+\setminus C_i\) has nonempty connected overlap with
\(U_0\).  Consequently the union
\[
 U_0\cup\bigcup_{i=1}^N\bigl(D_i^+\setminus C_i\bigr)
 =U\setminus\bigcup_{i=1}^N C_i
\]
is path connected.
\end{proof}